%%%%%%%%%%%%%%%%%%%%%%%%%%%%%%%%%%%%%%%%%%%%%%%%%%%%%%%%%%%%%%%%%%%%%%%%%%%%%%%%
%% 
\cleardoubleoddpage%  Make sure to start each chapter on a new odd page
\chapter{Tools under Investigation}
\label{sec:tools}


\section{Selection Criteria}
\label{sec:tools:criteria}

A big consideration of the tool selection was the availability and accessibility of the tools themselves. Programs like IDA Pro were excluded for the financial hurdle that they present. 
Preference was given to tools that are readily obtainable and usable in a standard analysis environment.

Static analysis tools that are used for managed code programs were also selected, simply because malware is not only written in languages that compile to machine code. 
This ensures that different types of compiled programs, including those targeting virtual machines such as the Java Virtual Machine or the .NET runtime, are represented in the evaluation.

The selected tools are widely used within the reverse engineering community. Established tools with active user bases and ongoing development were preferred, 
as they are more likely to be encountered in real world analysis scenarios.

Finally, the tools were selected based on their suitability for static analysis tasks. Each tool provides functionality for disassembly, decompilation, or both, 
and supports workflows relevant to malware analysis. Together, these criteria ensure that the selected tools provide a balanced and practical foundation 
for the comparative evaluation conducted in this thesis.

\section{Ghidra}
\label{sec:tools:ghidra}
Ghidra is, not just a tool, but a reverse engineering framework, developed and publicly released by the National Security Agency (NSA) in 2019 to support 
cybersecurity professionals, researchers, and educators. According to the NSA, the tool was designed to address complex software analysis challenges by providing a 
comprehensive suite of capabilities including disassembly, decompilation, debugging, and collaborative reverse engineering across multiple platforms. 
Its public release aimed to strengthen cybersecurity education, improve accessibility to advanced reverse engineering tools, and support malware analysis and vulnerability research. 
By making Ghidra freely available, the NSA positioned it as a significant resource for both professional analysts and academic environments.
Also, in their own words, they were trying to "level the playing field" for beginner analysts. The agency itself profits from open sourcing the framework by being able 
to hire new employees already familiar with a tool that they will actually then use in their work there. One could argue they also profit from additions and bugfixes from 
the open source community, as the project itself has a public github page.
\cite{bib:nsa-ghidra-release}

\section{Binary Ninja}
\label{sec:tools:binaryninja}
Binary Ninja is a static reverse engineering platform developed by Vector 35, a US-based company specialising in reverse engineering tools. 
The software originated from a private tool developed for use by a Capture the Flag (CTF) team. Vector 35 was subsequently founded to rebuild the tool as a commercial product 
\cite{bib:binaryninja_faq}. The original Python prototype was released under the GNU General Public License version 2, but it does not share source code with the current implementation 
\cite{bib:binaryninja_faq}.

The current version of Binary Ninja is proprietary software, although several components of the platform are available as open-source projects. These include its public API, 
debugger, and architecture definitions. Binary Ninja is offered in multiple editions, including a functionally restricted free edition and several paid editions intended for personal, 
commercial, and enterprise use \cite{bib:binaryninja_faq}. The desktop application supports Linux, macOS, and Windows.


\section{Cutter (Rizin GUI)}
\label{sec:tools:cutter}
Cutter is a free and open-source reverse engineering platform that provides a graphical user interface for the Rizin reverse engineering framework. 
It is implemented in C++ using the Qt application framework and is available for Linux, macOS, and Windows \cite{bib:cutter}. Rizin serves as Cutter's core analysis engine, 
while Cutter presents its functionality through a graphical and customizable user interface.

The project aims to make the extensive functionality of Rizin more accessible without requiring users to interact exclusively with its command-line interface. Nevertheless, 
Cutter also includes an integrated terminal through which Rizin commands can be executed directly. As both Cutter and Rizin are open-source projects, their source code can be inspected, 
modified, and extended by the community.

\section{Bytecode Viewer}
\label{sec:tools:bytecode-viewer}
Bytecode Viewer is a free and open-source Java and Android reverse-engineering suite released under the GNU General Public License version 3. It provides a graphical interface 
for inspecting and decompiling compiled Java classes and archives. Files can be opened through the application or imported by drag and drop, and supported formats include individual 
class files, JAR and WAR archives, as well as several Android formats \cite{bib:bytecode_viewer_github}.

Rather than relying on a single decompilation engine, Bytecode Viewer integrates several Java decompilers, including CFR, Procyon, FernFlower, Krakatau, JADX, and JD-GUI. 
It also provides bytecode disassemblers, search functionality, resource navigation, and configurable viewing panes that can display the output of up to three decompilers side 
by side \cite{bib:bytecode_viewer_github}. For the practical evaluation in this thesis, FernFlower was selected as the decompilation engine and used consistently with its default configuration.

Bytecode Viewer was included in this comparison to complement ILSpy with a tool designed for examining Java bytecode. Java applications and individual Java components may be 
encountered during malware analysis, making support for class files and JAR archives relevant to static analysis. Its integration of several decompilation and inspection components 
additionally distinguishes Bytecode Viewer from standalone decompilers that provide only reconstructed source code.

\section{ILSpy}
\label{sec:tools:ilspy}
ILSpy is a free and open-source .NET assembly browser and decompiler developed as part of the ICSharpCode project. It is distributed under the MIT License and 
provides a cross-platform desktop interface for Windows, Linux, and macOS \cite{bib:ilspy}. Its main purpose is the decompilation of .NET assemblies back into C\# source code. 
In addition to individual classes and methods, ILSpy can decompile entire projects and provides functionality for searching and navigating types, methods, and properties, 
as well as inspecting assembly metadata \cite{bib:ilspy}.

ILSpy was included in this comparison because .NET and C\# are also relevant in malware analysis. Malware samples may be written entirely in C\#, while others contain individual 
stages or components implemented as .NET assemblies. In such cases, a specialized .NET decompiler can provide significantly more useful high level output than a general purpose 
native-code reverse engineering tool.


%% 
%%%%%%%%%%%%%%%%%%%%%%%%%%%%%%%%%%%%%%%%%%%%%%%%%%%%%%%%%%%%%%%%%%%%%%%%%%%%%%%%
